%% ======================================================================
\section{PAC-Bayesian generalization bounds}
\label{chap:bounds}
%% ======================================================================

\begin{objectives}
\guiding{How much can the true risk of a randomly drawn voter exceed its
empirical risk, when the distribution of the voters has been chosen with the data?}
This section bounds the true risk of the Gibbs classifier $\GQ$. We first
explain why a useful bound must hold simultaneously for all posteriors, then
state and prove a general PAC-Bayesian theorem whose proof has only three steps
(Jensen, change of measure, Markov). The classical bounds of McAllester,
Seeger--Maurer and Catoni, and the PAC-Bayes-$\lambda$ bound, are all obtained
as corollaries, and we compare them numerically to learn how to read the
different terms of a bound.
\end{objectives}

\subsection{Why the bound must hold for all posteriors}

Before stating any PAC-Bayesian bound, we must understand what kind of
guarantee we need. A learning algorithm looks at the data before choosing its
output. If we compute a confidence bound for the chosen predictor as if it had
been fixed in advance, we forget this selection step, and the bound can be badly
wrong. The experiment reported in \cref{fig:uniformity} shows how badly. We
draw $1000$ classifiers that answer completely at random, so that each of them
has a true risk of exactly $0.5$, and we evaluate them on a sample of $m=100$
examples. By chance alone, the best of them has an empirical risk of about
$0.34$ on average. If we now apply Hoeffding's inequality to this best
classifier as if it were fixed, we obtain an ``upper bound'' which is below
$0.5$ in $99.8\%$ of the runs, although it should fail in at most $5\%$ of
them. The Occam bound of \cref{sec:occam}, which pays $\ln 1000$ for the
selection, fails in only $0.3\%$ of the runs.

\begin{figure}[t]
\centering
\includegraphics[width=0.72\linewidth]{uniformity}
\caption{Selection bias. Distribution, over $2000$ repetitions, of the smallest
empirical risk among $1000$ random classifiers of true risk $0.5$, measured on
$m=100$ examples. A confidence bound computed for the selected classifier as
if it had been fixed in advance is violated in almost every run.}
\label{fig:uniformity}
\end{figure}

This is why all the bounds of this section are \emph{uniform}: they hold, with
probability at least $1-\delta$, for all posteriors $Q$ at the same time. The
posterior can then be chosen by any procedure that looks at the data ---
including the minimization of the bound itself --- without invalidating the
guarantee. The price of uniformity is the complexity term $\KL(Q\|P)$, which
plays the role that $\ln|\Hcal|$ played in the Occam bound. The next subsection
shows that this price is enough, with a single theorem.

\subsection{A general PAC-Bayesian theorem}

Rather than proving each classical bound separately, we prove one theorem that
contains them all, in the form given by \citet{germain2009pac} and
\citet{begin2016pac}. It is parametrized by a convex function $\Delta$, which
measures the discrepancy between an empirical risk and a true risk.

\begin{theorem}{General PAC-Bayesian theorem \citep{germain2009pac,begin2016pac}}{general}
Let $\Delta:[0,1]\times[0,1]\to\R$ be a \emph{jointly convex} function. For any
distribution $\D$ on $\X\times\Y$, any set of voters $\Hcal$, any loss
$\ell$ with values in $[0,1]$, any prior $P\in\M(\Hcal)$ and any
$\delta\in(0,1]$, with probability at least $1-\delta$ over $S\sim\D^m$,
\begin{equation}\label{eq:general}
\forall Q\in\M(\Hcal),\quad
\Delta\big(\RS(\GQ),\RD(\GQ)\big)\leq\frac1m\left[\KL(Q\|P)+\ln\frac{\mathcal I_\Delta(m)}{\delta}\right],
\end{equation}
where $\mathcal I_\Delta(m)=\Esp{S'\sim\D^m}{\Esp{h\sim P}{e^{m\Delta(\widehat R_{S'}(h),\RD(h))}}}$.
\end{theorem}

In words, with high probability, and for every posterior at once, the gap
between the empirical and the true Gibbs risk, measured by $\Delta$, is at most
a complexity per example: the KL divergence plus a confidence term, divided by
$m$. Only $\mathcal I_\Delta(m)$ is left to compute, and it involves the prior
alone.

Before the proof, let us describe its strategy. We want to control a
quantity that involves the posterior $Q$, which depends on the data, and we only
know how to control quantities that involve voters fixed in advance. The proof
therefore moves the problem in three steps. First, it replaces the discrepancy of
the \emph{average} voter by the average discrepancy of the voters (Jensen).
Second, it replaces the average under $Q$ by an average under the prior $P$, at
the price $\KL(Q\|P)$ (change of measure). Third, it controls the resulting
quantity, which no longer involves $Q$, on average over the samples (Markov).
The function $\Delta$ only matters in the last step, and each choice of
$\Delta$ will give one of the classical bounds.

\begin{proof}
Fix a sample $S$ and a posterior $Q$. The first step uses Jensen's
inequality. Since $\Delta$ is jointly convex and the two risks of $\GQ$ are
$Q$-expectations of the risks of the voters,
\begin{align*}
m\,\Delta\big(\RS(\GQ),\RD(\GQ)\big)
&=m\,\Delta\Big(\Esp{h\sim Q}{\RS(h)},\Esp{h\sim Q}{\RD(h)}\Big)\\
&\leq\Esp{h\sim Q}{m\,\Delta\big(\RS(h),\RD(h)\big)}.
\end{align*}
The second step is the change of measure: Lemma~\ref{lem:dv} with
$\phi(h)=m\,\Delta(\RS(h),\RD(h))$ gives
\[
\Esp{h\sim Q}{m\,\Delta\big(\RS(h),\RD(h)\big)}\leq\KL(Q\|P)+\ln\Esp{h\sim P}{e^{m\Delta(\RS(h),\RD(h))}}.
\]
The posterior has now disappeared from the last term, which only involves the
prior; it is a random variable $Z(S)=\E_{h\sim P}\,e^{m\Delta(\RS(h),\RD(h))}$
that depends on the sample only. The third step bounds this variable. It is
non-negative, so Markov's inequality ensures that, with probability at least
$1-\delta$ over $S$, $Z(S)\leq\frac1\delta\E_{S'}Z(S')$. Since $P$ does not
depend on $S'$, Fubini's theorem allows us to exchange the two expectations,
$\E_{S'}Z(S')=\E_{h\sim P}\E_{S'}e^{m\Delta(\widehat R_{S'}(h),\RD(h))}=\mathcal I_\Delta(m)$.

On the event of probability at least $1-\delta$ given by the third step,
inequality \eqref{eq:general} holds for \emph{every} $Q$, because the first
two steps are deterministic and valid for all posteriors simultaneously.
\end{proof}

The structure of this proof deserves attention, because it explains most of
what follows. The only probabilistic statement is made in the third step, and
it only concerns the prior: this is where the prior must be independent of the
sample, since we need to exchange $\E_{S'}$ and $\E_{h\sim P}$ with $P$
fixed. The posterior, on the other hand, only appears in the deterministic
steps, and it may depend on $S$ in an arbitrary way. This asymmetry is what
makes PAC-Bayesian bounds usable as learning objectives: we are free to
minimize them over $Q$.

To obtain a usable bound, it remains to choose $\Delta$ so that $\mathcal
I_\Delta(m)$ can be computed or bounded. For a \emph{fixed} voter $h$, $m\RS(h)$
is a sum of $m$ independent variables in $[0,1]$ with mean $\RD(h)$, so the
concentration results of \cref{chap:tools} apply. Each of them yields one of the
classical bounds, and the next subsection goes through them in turn: the kl form
(the tightest), its relaxation by Pinsker's inequality (the most readable),
Catoni's bound, and two linear bounds (the easiest to optimize).

\subsection{The classical bounds}

\subsubsection{Seeger's bound (kl form)}

The most natural choice is the binary kl itself, since Maurer's lemma tells us
exactly how large its exponential moment is.

\begin{corollary}{\citet{seeger2002pac,langford2005tutorial}, with the constant of \citet{maurer2004note}}{seeger}
For $m\geq8$, with probability at least $1-\delta$ over $S\sim\D^m$, for all $Q\in\M(\Hcal)$,
\begin{equation}\label{eq:seeger}
\kl\big(\RS(\GQ)\,\|\,\RD(\GQ)\big)\leq\frac{\KL(Q\|P)+\ln\frac{2\sqrt m}{\delta}}{m},
\end{equation}
that is, $\RD(\GQ)\leq\klup\Big(\RS(\GQ),\ \frac{\KL(Q\|P)+\ln\frac{2\sqrt m}{\delta}}{m}\Big)$.
\end{corollary}
\begin{proof}
Take $\Delta=\kl$, which is jointly convex. By Maurer's lemma
(Lemma~\ref{lem:chernoff}), for every fixed $h$,
$\E_{S'}e^{m\kl(\widehat R_{S'}(h)\|\RD(h))}\leq\xi(m)\leq2\sqrt m$, hence
$\mathcal I_{\kl}(m)\leq2\sqrt m$.
\end{proof}

In plain words, the true Gibbs risk is at most the largest value that is still
``plausible'', in the sense of the kl, given the observed Gibbs risk and a
budget equal to the complexity per example. Seeger's bound has no free
parameter and no closed form, but it is computed in a few microseconds by the
bisection of \cref{algo:klinv}. As the numerical comparison below will show, it
is among the tightest bounds available, and it is the one we use by default to
compute risk certificates. It inherits the
fast rate of the kl: when the empirical risk is zero, it gives
$\RD(\GQ)\leq1-e^{-\varepsilon}\leq\varepsilon$ with
$\varepsilon=\frac1m[\KL(Q\|P)+\ln\frac{2\sqrt m}\delta]$, a rate in $1/m$.
Its only weakness is that it is implicit, which makes it hard to read; the next
bound trades some tightness for readability.

\subsubsection{McAllester's bound}

Relaxing the kl with Pinsker's inequality gives the historical PAC-Bayesian
bound, whose form is the easiest to interpret.

\begin{corollary}{\citet{mcallester1999some,mcallester2003pac}, in the form of \citet{maurer2004note}}{mcallester}
For $m\geq 8$, with probability at least $1-\delta$ over $S\sim\D^m$, for all $Q\in\M(\Hcal)$,
\begin{equation}\label{eq:mcallester}
\RD(\GQ)\leq\RS(\GQ)+\sqrt{\frac{\KL(Q\|P)+\ln\frac{2\sqrt m}{\delta}}{2m}}.
\end{equation}
\end{corollary}
\begin{proof}
Apply Pinsker's inequality $\kl(q\|p)\geq2(q-p)^2$ to \eqref{eq:seeger}.
\end{proof}

McAllester's bound reads ``true risk $\leq$ empirical risk $+$ complexity'',
where the complexity is a square root of the KL per example. It has the same
structure as the VC or Rademacher bounds of the first part of the course, the
capacity of the class being replaced by the divergence of the chosen posterior
from the prior. Its drawback is the $1/\sqrt m$ rate, which it keeps even when
the empirical risk is zero; Seeger's bound is always at least as tight.

\begin{history}
The original bounds of \citet{mcallester1998some,mcallester1999some} had
slightly worse constants, with a numerator of the form
$\KL(Q\|P)+\ln\frac1\delta+O(\ln m)$ and a denominator $2m-1$. The kl form was
introduced by \citet{seeger2002pac}, see also \citet{langford2005tutorial}, and the clean
constant $\ln\frac{2\sqrt m}{\delta}$, valid for any loss in $[0,1]$, is due to
\citet{maurer2004note}.
\end{history}

\subsubsection{Catoni's bound}

Can we avoid the $\ln 2\sqrt m$ term altogether? \citet{catoni2007pac} uses a
different comparison function, for which the exponential moment is not bounded
by $2\sqrt m$ but by $1$, at the cost of a parameter that must be fixed in
advance.

\begin{corollary}{\citet{catoni2007pac}, see also \citet{germain2009pac}}{catoni}
For any fixed $C>0$, with probability at least $1-\delta$ over $S\sim\D^m$, for all $Q$,
\begin{equation}\label{eq:catoni}
\RD(\GQ)\leq\frac{1}{1-e^{-C}}\left[1-\exp\Big(-C\,\RS(\GQ)-\frac{\KL(Q\|P)+\ln\frac1\delta}{m}\Big)\right].
\end{equation}
\end{corollary}
\begin{proof}
Take $\Delta(q,p)=\mathcal F_C(p)-C\,q$ with $\mathcal F_C(p)=-\ln\big(1-p(1-e^{-C})\big)$,
which is convex. For a fixed $h$ with risk $p=\RD(h)$, the convexity of
$u\mapsto e^{-Cu}$ on $[0,1]$ gives $e^{-Cu}\leq1-u(1-e^{-C})$, hence, by
independence of the examples, $\E_{S'}e^{-Cm\widehat R_{S'}(h)}\leq\big(1-p(1-e^{-C})\big)^m=e^{-m\mathcal F_C(p)}$.
Therefore $\E_{S'}e^{m\Delta(\widehat R_{S'}(h),p)}\leq1$ and $\mathcal I_\Delta(m)\leq1$.
Rearranging $\mathcal F_C(\RD(\GQ))\leq C\RS(\GQ)+\frac1m[\KL+\ln\frac1\delta]$ gives \eqref{eq:catoni}.
\end{proof}

The parameter $C$ weighs the empirical risk against the complexity. It cannot be
optimized after seeing the data for free, but the price is small. If we apply
the bound with confidence $\delta/k$ to each value of a grid $C_1,\dots,C_k$
fixed in advance, a union bound shows that it holds for all of them
simultaneously; the best value can then be chosen at a cost of only $\ln k$ in
the numerator. With such a grid, Catoni's bound is numerically very
close to Seeger's, and even marginally smaller when the KL is small, since it
does not pay the $\ln 2\sqrt m$ term. Using $1-e^{-x}\leq x$, it also implies
the \emph{linear} bound
$\RD(\GQ)\leq\frac{1}{1-e^{-C}}\big[C\,\RS(\GQ)+\frac{\KL(Q\|P)+\ln(1/\delta)}{m}\big]$.

\subsubsection{Linear bounds and the PAC-Bayes-\texorpdfstring{$\lambda$}{lambda} bound}

The last two bounds of this section are \emph{linear} in $\RS(\GQ)$ and in
$\KL(Q\|P)$. This is a weakness as far as tightness is concerned, but a great
strength for learning: minimizing a linear combination of the empirical risk and
of the KL over $Q$ has an explicit solution, the Gibbs posterior of
\cref{sec:gibbs_posterior}.

\begin{corollary}{Hoeffding-based bound \citep{catoni2007pac,alquier2024user}}{hoeffding_pb}
For any fixed $\lambda>0$, with probability at least $1-\delta$ over $S\sim\D^m$, for all $Q$,
\begin{equation}\label{eq:catoni_hoeffding}
\RD(\GQ)\leq\RS(\GQ)+\frac{\KL(Q\|P)+\ln\frac1\delta}{\lambda m}+\frac{\lambda}{8}.
\end{equation}
\end{corollary}
\begin{proof}
Take $\Delta(q,p)=\lambda(p-q)$, which is linear, hence convex. By Hoeffding's
lemma (Lemma~\ref{lem:hoeffding}) applied with parameter $\lambda m$,
$\E_{S'}e^{\lambda m(\RD(h)-\widehat R_{S'}(h))}\leq e^{\lambda^2m/8}$, so that
$\frac1m\ln\mathcal I_\Delta(m)\leq\lambda^2/8$. Divide by $\lambda$.
\end{proof}

The bound reads ``empirical risk $+$ complexity$/\lambda$ $+$ $\lambda/8$'': a
large $\lambda$ makes the KL cheap but increases the last term, and a small
$\lambda$ does the opposite. The Hoeffding-based bound is the favourite of the
statistical literature on PAC-Bayes \citep{alquier2024user}, because it leads to clean oracle inequalities.
Its parameter $\lambda$ must again be fixed in advance. The next bound, obtained
by relaxing Seeger's bound in a clever way, does not have this limitation.

\begin{corollary}{PAC-Bayes-$\lambda$ bound \citep{thiemann2017strongly}}{lambda}
With probability at least $1-\delta$ over $S\sim\D^m$, for all $Q\in\M(\Hcal)$ and all
$\lambda\in(0,2)$ \emph{simultaneously},
\begin{equation}\label{eq:lambda}
\RD(\GQ)\leq\frac{\RS(\GQ)}{1-\frac\lambda2}+\frac{\KL(Q\|P)+\ln\frac{2\sqrt m}{\delta}}{\lambda\big(1-\frac\lambda2\big)m}.
\end{equation}
\end{corollary}
\begin{proof}
Let $\hat q=\RS(\GQ)$, $p=\RD(\GQ)$ and $\varepsilon=\frac1m[\KL(Q\|P)+\ln\frac{2\sqrt m}\delta]$.
On the event of Corollary~\ref{cor:seeger}, if $p>\hat q$, the refined Pinsker inequality gives
$(p-\hat q)^2\leq2p\,\varepsilon$, that is $p-\hat q\leq\sqrt{2p\varepsilon}$.
For any $\lambda>0$, $\sqrt{2p\varepsilon}=\sqrt{\lambda p\cdot\frac{2\varepsilon}{\lambda}}\leq\frac{\lambda p}{2}+\frac{\varepsilon}{\lambda}$
by the inequality between arithmetic and geometric means.
Hence $p(1-\frac\lambda2)\leq\hat q+\frac\varepsilon\lambda$. Since this is a
deterministic consequence of \eqref{eq:seeger}, it holds for all $\lambda$ at once.
\end{proof}

The PAC-Bayes-$\lambda$ bound keeps the linear structure of the previous one,
but it is a deterministic consequence of Seeger's bound. Because it holds for
all $\lambda$ simultaneously, $\lambda$ can be optimized freely, at no cost. For a fixed $Q$, the optimal value is
\begin{equation}\label{eq:lambda_star}
\lambda^\star=\frac{2}{\sqrt{\frac{2m\RS(\GQ)}{\KL(Q\|P)+\ln\frac{2\sqrt m}\delta}+1}+1},
\end{equation}
and for this value the PAC-Bayes-$\lambda$ bound is smaller than the
relaxation $\hat q+\sqrt{2\hat q\varepsilon}+2\varepsilon$ of Seeger's bound
given by Lemma~\ref{lem:pinsker}, but larger than Seeger's bound itself, from
which it is derived. \citet{thiemann2017strongly} show moreover that
\eqref{eq:lambda} is jointly quasiconvex in $(Q,\lambda)$ under mild conditions,
which is the basis of the alternating minimization algorithm of \cref{sec:algo_fo}.

\subsection{Reading and comparing the bounds}

We now have five bounds; to choose among them, we must learn to read them.
All these bounds share the same anatomy. The unknown true risk is bounded by a
function of two quantities: the empirical risk of the posterior, which measures
how well it fits the data, and the complexity per example
$\frac1m[\KL(Q\|P)+\ln\frac{2\sqrt m}\delta]$, which measures how far the
posterior has moved away from the prior. The first term decreases when $Q$
concentrates on good voters, the second increases at the same time, and the
bounds only differ in the way these two ingredients are combined. The
confidence term is almost negligible: dividing $\delta$ by ten only adds
$\ln 10\approx2.3$ nats to the numerator.

\begin{figure}[t]
\centering
\includegraphics[width=0.75\linewidth]{bound_anatomy}
\caption{Decomposition of McAllester's and Seeger's bounds into the empirical
term (grey) and the complexity term, for $m=1000$, $\delta=0.05$ and
$\RS(\GQ)=0.12$, as the KL divergence grows.}
\label{fig:bound_anatomy}
\end{figure}

\Cref{fig:bound_anatomy} and \cref{tab:bound_anatomy} make this concrete. We fix
$m=1000$, $\delta=0.05$ and an empirical Gibbs risk of $0.12$, and we let the KL
divergence grow from $0$ to $50$ nats. With $\KL=0$, the complexity per example
is only $0.007$, and all bounds lie between $0.16$ and $0.18$. The complexity
term of McAllester's bound is already $0.06$, half the empirical risk, while
Seeger's bound adds only $0.04$: the square root of Pinsker's inequality is
expensive when the risk is far from $1/2$.

As the KL grows, all bounds
deteriorate, but slowly: $50$ nats, which is the cost of a posterior putting
all its mass on one voter among $e^{50}\approx5\times10^{21}$, still leaves a
non-trivial certificate of $0.26$. Catoni's bound optimized on a grid and
Seeger's bound are almost indistinguishable, the former being slightly smaller
for small KL. The PAC-Bayes-$\lambda$ bound, finally, is tighter than
McAllester's bound when the KL is small but not when it is large.

\begin{table}[t]
\centering\small
\caption{Values of the bounds for $\RS(\GQ)=0.12$, $m=1000$, $\delta=0.05$.
Catoni's bound is optimized on a grid of $30$ values of $C$ with a union bound.}
\label{tab:bound_anatomy}
\begin{tabular}{rccccc}
\toprule
$\KL(Q\|P)$ & $\frac1m[\KL+\ln\frac{2\sqrt m}{\delta}]$ & McAllester & PAC-Bayes-$\lambda$ & Catoni (grid) & Seeger \\
\midrule
0 & 0.0071 & 0.180 & 0.169 & 0.160 & 0.162 \\
2 & 0.0091 & 0.188 & 0.177 & 0.166 & 0.168 \\
5 & 0.0121 & 0.198 & 0.187 & 0.175 & 0.177 \\
10 & 0.0171 & 0.213 & 0.204 & 0.187 & 0.189 \\
20 & 0.0271 & 0.236 & 0.232 & 0.207 & 0.209 \\
50 & 0.0571 & 0.289 & 0.307 & 0.256 & 0.256 \\
\bottomrule
\end{tabular}
\end{table}

\Cref{fig:bounds} shows the same comparison from two other angles. As a
function of the sample size, the difference between the kl-based bounds and
McAllester's bound is most visible when the empirical risk is zero: the former
decrease like $1/m$, the latter like $1/\sqrt m$. As a function of the empirical
risk, the bounds coincide near $1/2$ and separate when the risk is small, which
is precisely the regime of good classifiers.

\begin{figure}[t]
\centering
\includegraphics[width=\linewidth]{bounds_comparison}
\caption{Comparison of the bounds for $\delta=0.05$. Left: as a function of
$m$, for $\RS(\GQ)=0.1$ and $\KL(Q\|P)=5$. Middle: same with $\RS(\GQ)=0$: the
kl-based bounds (Seeger, and Catoni and PAC-Bayes-$\lambda$ which approximate it)
decrease in $1/m$ whereas McAllester's bound decreases in $1/\sqrt m$.
Right: as a function of the empirical risk for $m=1000$ and $\KL=10$.}
\label{fig:bounds}
\end{figure}

The following table summarizes the section. The first column gives the
comparison function used in Theorem~\ref{thm:general}, and the second the bound
on its exponential moment that makes the theorem explicit.

\begin{center}
\renewcommand{\arraystretch}{1.3}
\small
\begin{tabular}{lccp{4.2cm}}
\toprule
Bound & $\Delta(q,p)$ & $\ln\mathcal I_\Delta(m)\leq$ & Main features \\
\midrule
Seeger \eqref{eq:seeger} & $\kl(q\|p)$ & $\ln 2\sqrt m$ & among the tightest; implicit \\
McAllester \eqref{eq:mcallester} & $2(q-p)^2$ & $\ln 2\sqrt m$ & readable; slow rate $1/\sqrt m$ \\
Catoni \eqref{eq:catoni} & $\mathcal F_C(p)-Cq$ & $0$ & tight on a grid of $C$ \\
Hoeffding \eqref{eq:catoni_hoeffding} & $\lambda(p-q)$ & $\lambda^2m/8$ & linear; $\lambda$ fixed in advance \\
PAC-Bayes-$\lambda$ \eqref{eq:lambda} & (from Seeger) & $\ln 2\sqrt m$ & linear, all $\lambda$ at once \\
\bottomrule
\end{tabular}
\end{center}

\begin{example}[Three posteriors on the same ensemble]
Consider an ensemble of $n=100$ voters with the uniform prior $P$, $m=1000$ training
examples and $\delta=0.05$, so that $\ln\frac{2\sqrt m}{\delta}=\ln(1264.9)\approx7.14$.
Suppose first that $Q$ is uniform over the $10$ voters with the smallest
empirical risk, and that its empirical Gibbs risk is $0.12$. The divergence is
$\KL(Q\|P)=\ln\frac{100}{10}\approx2.30$, so that
$\varepsilon=\frac{2.30+7.14}{1000}\approx0.0094$. McAllester's bound gives
$0.12+\sqrt{0.0094/2}\approx0.189$, the PAC-Bayes-$\lambda$ bound (with
$\lambda^\star\approx0.33$) gives $0.178$, and Seeger's bound gives
$\klup(0.12,0.0094)\approx0.169$. If instead we keep $Q=P$, the KL vanishes but
the Gibbs risk is larger, say $0.20$, and Seeger's bound gives $0.251$. At the
other extreme, the Dirac on the best voter, with empirical risk $0.10$, costs
$\KL=\ln100\approx4.61$ and Seeger's bound gives $0.152$: this is the Occam bound.

In this example, the best \emph{Gibbs} certificate is thus obtained by the
Dirac, because $\ln100$ nats is a small price for $m=1000$ examples. But the
Gibbs risk is not the end of the story: the \emph{majority vote} of the $10$ best
voters can be much more accurate than any of them, and certifying it requires
spread-out posteriors, as we will see in \cref{chap:mv}. In general, the best
trade-off between empirical risk and KL is given by the Gibbs posterior of
\cref{sec:gibbs_posterior}.
\end{example}

\begin{running}
Let us certify the Gibbs classifier of our forest, that is a tree drawn at
random, with the uniform prior and the uniform posterior, so that $\KL(Q\|P)=0$.
We need examples that the trees have not seen; for now, let us admit that each
tree may be evaluated on the examples left out of its bootstrap sample (its
\emph{out-of-bag} examples), a construction justified in \cref{sec:oob}. The
average out-of-bag error of the trees is $0.174$, computed with at least $747$
examples per tree. With $\delta=0.05$, the complexity per example is
$\ln(2\sqrt{747}/0.05)/747=0.0094$, McAllester's bound gives $0.243$ and Seeger's
bound $0.230$. The test Gibbs risk, which the bound did not see, is $0.190$: the
certificate is valid and only four points above. The difficulty is elsewhere:
the vote has a test risk of $0.044$, and nothing in this section says anything
about it. This is the object of \cref{chap:mv}.
\end{running}

\begin{remark}[(Bernstein-type bounds)]
When the variance of the loss is small, the kl bound can be further improved
by PAC-Bayes-Bernstein inequalities, which involve an empirical variance term
\citep{tolstikhin2013pac,mhammedi2019pac}. The same idea is used for majority
votes in \cref{sec:cctnd} \citep{wu2021chebyshev}.
\end{remark}

\begin{keypoints}
A generalization bound for a data-dependent predictor must be uniform, otherwise
the selection bias makes it meaningless. All the classical PAC-Bayesian bounds
follow from Theorem~\ref{thm:general}, whose only probabilistic step uses the
independence between the prior and the sample; the posterior can therefore
depend on the data arbitrarily. The complexity term is $\KL(Q\|P)/m$: a posterior
close to the prior is cheap. Seeger's kl bound and Catoni's bound optimized on a
grid are the tightest, McAllester's bound is the easiest to read, and linear
bounds such as PAC-Bayes-$\lambda$ are the easiest to minimize.
\end{keypoints}

\begin{quickchecks}
\item Which step of the proof of Theorem~\ref{thm:general} requires the prior to be independent of the sample?
\item Why may the posterior depend on the data without invalidating the bound?
\item Which bound has a rate in $1/m$ when the empirical risk is zero, McAllester's or Seeger's?
\item Why can $\lambda$ be optimized freely in the PAC-Bayes-$\lambda$ bound, but not in the Hoeffding-based bound?
\end{quickchecks}

\begin{exercises}
\begin{exercise}
Starting from Theorem~\ref{thm:general} with $\Delta(q,p)=2(q-p)^2$, prove McAllester's
bound directly, without going through the kl.
\end{exercise}
\begin{exercise}
For $m=500$, $\delta=0.05$, $\RS(\GQ)=0.05$ and $\KL(Q\|P)\in\{0,5,20\}$, compute
McAllester's and Seeger's bounds, and compare with \cref{tab:bound_anatomy}.
\end{exercise}
\begin{exercise}
Using the linear form of Catoni's bound given after Corollary~\ref{cor:catoni}, find
the value of $C$ that minimizes it when $\RS(\GQ)=0.12$, $\KL(Q\|P)=5$, $m=1000$ and
$\delta=0.05$, and compare the result with \cref{tab:bound_anatomy}.
\end{exercise}
\begin{exercise}
Reproduce the experiment of \cref{fig:uniformity} for $m\in\{100,1000,10000\}$.
Show that the failure frequency of the naive bound does not go to zero when $m$ grows, and
explain why by comparing $\sqrt{\ln(1/\delta)/(2m)}$ with the typical deviation of the
minimum of $1000$ empirical risks.
\end{exercise}
\begin{exercise}
\textbf{(Comprehension)} In the proof of Theorem~\ref{thm:general}, which step uses the
joint convexity of $\Delta$, which step uses the independence of the prior from the
sample, and which step is the only probabilistic one? Why does the bound hold for all
posteriors although Markov's inequality is applied only once?
\end{exercise}
\begin{exercise}
\textbf{(Comprehension)} True or false? Justify briefly. (a)~For the same
$\RS(\GQ)$, $\KL(Q\|P)$, $m$ and $\delta$, Seeger's bound is never larger than
McAllester's bound. (b)~Doubling the sample size $m$ halves the complexity term of
McAllester's bound. (c)~In Catoni's bound, the parameter $C$ may be chosen after
seeing the data, since the bound holds for all posteriors. (d)~Dividing $\delta$ by
$10$ adds about $2.3$ nats to the numerator of the complexity term.
\end{exercise}
\begin{exercise}
\textbf{(Application)} An ensemble of $n=50$ voters has a uniform prior. With $m=500$
examples and $\delta=0.05$, the posterior $Q$ uniform over the $5$ best voters has an
empirical Gibbs risk of $0.08$. Compute $\KL(Q\|P)$ and the complexity per example
$\varepsilon=\frac1m[\KL(Q\|P)+\ln\frac{2\sqrt m}{\delta}]$, then McAllester's bound, the
PAC-Bayes-$\lambda$ bound at $\lambda^\star$, Seeger's bound (with \texttt{pbtools}), and
Catoni's bound for $C\in\{0.5,1,2\}$. Which value of $C$ is legitimate to report?
\end{exercise}
\begin{exercise}
\textbf{(Application)} A posterior has zero empirical Gibbs risk on $m=200$ examples,
with $\KL(Q\|P)=3$ and $\delta=0.05$. Compute Seeger's bound (using its closed form
$1-e^{-\varepsilon}$ when $\RS(\GQ)=0$) and McAllester's bound. How many examples does
each of them need to certify a Gibbs risk below $0.05$ (with the same KL)?
\end{exercise}
\begin{exercise}
\textbf{(Proof)} Show that the binary kl is \emph{jointly} convex in $(q,p)\in[0,1]\times(0,1)$,
as required by Theorem~\ref{thm:general}. \emph{Hint:} for $a>0$, $b>0$, write
$a\ln\frac ab=a\,f\big(\frac ba\big)$ with the convex function $f(u)=-\ln u$, and show that
the perspective $(a,b)\mapsto a\,f(b/a)$ of a convex function is jointly convex.
\end{exercise}
\begin{exercise}
\textbf{(Proof)} Show that, for fixed $\hat q>0$ and $\varepsilon>0$, the function
$\lambda\mapsto\frac{\hat q}{1-\lambda/2}+\frac{\varepsilon}{\lambda(1-\lambda/2)}$ is minimized on $(0,2)$
at $\lambda^\star=2/\big(\sqrt{1+2\hat q/\varepsilon}+1\big)$, which is \eqref{eq:lambda_star}, and that
its minimal value is $\big(\sqrt{\hat q+\varepsilon/2}+\sqrt{\varepsilon/2}\big)^2$.
\end{exercise}
\end{exercises}
